\section{曲线与六维 $\cN{=}(2,0)$ 理论}\label{sec:6d}\index{SixDimensionalTheory@六维 $\cN{=}(2,0)$ 理论|hyperpage}
我们已经看到，纯 $\SU(2)$ 规范理论与含有一个味的 $\SU(2)$ 规范理论，其低能动力学都可以用复曲线 \eqref{curve_pure}、\eqref{curve_nf_1} 表示。本节的目的是说明，这些二维空间可以赋予物理解释。

本节将介绍的思想，最初是利用弦论和 M 理论的各种深刻性质得到的，包括 Calabi-Yau 紧致化、膜构造以及弦对偶。利用 Calabi-Yau 紧致化的方法可以追溯到文献 \cite{Kachru:1995wm,Kachru:1995fv,Klemm:1996bj}，而膜构造方法则由文献 \cite{Witten:1997sc} 引入。学习这些构造无疑有助于理解 $\cN{=}2$ 超对称动力学，反过来也一样。不过，本讲义并不讲授这些内容。

这里的论述更多依靠类比，而非严密的逻辑推演；作者有意选择这样的表述方式，使阅读所需的弦论和 M 理论知识尽可能少。对更多细节感兴趣的读者，应参阅原始论文，或文献 \cite{Lerche:1996xu,Giveon:1998sr} 这样的综述，以及文献 \cite{Gaiotto:2009hg} 的第 3 节。

\subsection{张力可变的弦}\index{VariableTensionString@张力可变的弦|hyperpage}\index{BPSEquation@BPS 方程|hyperpage}
回忆纯理论中电荷为 $n$、磁荷为 $m$ 的粒子的 BPS 质量公式，
\begin{equation}
M\ge |na+ma_D|=|\int_L\lambda|\label{bpsagain}
\end{equation}
其中 $L$ 是曲线上沿 $A$ 方向绕行 $n$ 次、沿 $B$ 方向绕行 $m$ 次的闭链。我们采用的基本思想，是认真看待这个方程：将四维粒子视为缠绕在闭链 $L$ 上的一条弦所产生的对象。此时，$\lambda$ 就类似于弦的张力。在本节中，我们在 $\lambda$ 的定义中引入一个 $2\pi i$ 因子，以简化公式。

为了使这一思想更具体，假设有一个以弦为激发的六维理论，\footnote{一些刚开始学习弦论的年轻读者，此时可能会疑惑：具有相对论性洛伦兹不变性的弦论，不是只有在玻色情形下的 $26$ 维、或超对称情形下的 $10$ 维中才可能存在吗？关键在于，教科书中的标准论证假定弦之间的相互作用是微扰性的。} 并假设这一理论定义在二维空间 $C$ 与四维闵可夫斯基空间 $\bR^{1,3}$ 的乘积上。再假设弦的张力依赖于这些额外维度的方向。也就是说，假定存在一个局部全纯的一形式 $\lambda$，使得以 $s$ 参数化的无穷小弦段的张力为
\begin{equation}
|\lambda|:=|\frac{\lambda}{ds}|ds,
\end{equation}
见图 \ref{fig:var_tension_string} 左侧。为明确起见，图中的二维空间取为环面。
\begin{figure}[h]
\[
\inc{variabletension}
\]
\caption{张力可变的弦。其张力由 $f(s)=\int^s_{P_0}\lambda$ 控制。\label{fig:var_tension_string}}
\end{figure}

从未紧致化的四维空间看，弦表现为一个粒子，其质量由可变张力的积分给出：
\begin{equation}
M=\int_L|\lambda|\label{m-l}
\end{equation}
利用三角不等式的积分形式，可以给出右端的下界：
\begin{equation}
\int_L|\lambda|\ge|\int_L\lambda|.\label{tri-ineq}
\end{equation}
考虑复平面中由
\begin{equation}
f(s)=\int^s_{P_0}\lambda
\end{equation}
定义、以 $s$ 为参数的曲线，就可以将这一不等式直观地呈现出来，其中 $P_0$ 是闭链 $L$ 上的一个固定点。式 \eqref{tri-ineq} 左端是参数曲线 $f(s)$ 的长度，而右端是 $f(s)$ 两个端点之间的距离，见图 \ref{fig:var_tension_string} 右侧。显然，前者长于后者，只有当曲线 $f(s)$ 本身是一条直线时才取等号。等价地，
\begin{equation}
\Arg\frac{\lambda}{ds}=\text{常数}.\label{curveBPSeq}
\end{equation}

当闭链 $L$ 在拓扑上平凡时，函数 $f(s)$ 的像本身也是一个闭环，式 \eqref{tri-ineq} 的右端为零。当闭链 $L$ 非平凡时，函数 $f(s)$ 的像可以是一条开线段。由于 $\lambda$ 全纯，线段两端的差只依赖于闭链 $L$ 的拓扑。设 $L$ 绕 $A$ 闭链 $n$ 次，绕 $B$ 闭链 $m$ 次。结合式 \eqref{m-l} 和 \eqref{tri-ineq}，得到
\begin{equation}
M\ge|na+ma_D|
\end{equation}
其中 $a$、$a_D$ 由关系
\begin{equation}
a=\int_A\lambda,\qquad
a_D=\int_B\lambda.
\end{equation}
定义。这就重现了 BPS 质量公式 \eqref{bpsagain}。此外，我们还知道，只有满足式 \eqref{curveBPSeq} 时，不等式才取等号。因而，将式 \eqref{curveBPSeq} 视为弦激发的 BPS 方程。

\subsection{由膜产生的张力可变的弦}
\subsubsection{一般思想}\index{MTheory@M 理论|hyperpage}\index{MTwoBrane@M2 膜|hyperpage}\index{MFiveBrane@M5 膜|hyperpage}
\begin{figure}[h]
\[
\inc{M5M2}
\]
\caption{张力可变的弦如何从更高维度产生。\label{fig:membrane-to-string}}
\end{figure}
有人或许会觉得，张力可变的弦有些奇怪。自然实现这种张力变化的一种方式，是将二维额外空间 $C$ 进一步嵌入四维环境空间 $X$ 中，并让覆盖 $C$ 的 $\Sigma$ 的两个叶片在 $X$ 的额外方向上彼此分离。再进一步假设，存在一个沿两个空间方向和一个时间方向延伸的膜，其端点可以落在 $\Sigma$ 的叶片上。图 \ref{fig:membrane-to-string} 描绘了这一情形。令 $z$ 为 $C$ 的坐标，$X$ 的复坐标为 $(z,x)$。于是，$\Sigma$ 的两个叶片定义了两个函数 $x_1(z)$ 和 $x_2(z)$。一个张力恒为 $|dx|\wedge|d\log z|$、悬置于两个叶片之间的膜，就可以视为一条张力可变的弦，其在给定 $z$ 值处的张力满足
\begin{equation}
\text{（$z$ 处的张力）}\ge\left|\int_{x_2(z)}^{x_1(z)}dx\wedge d\log z\right|
=\left|x_1\frac{dz}z-x_2\frac{dz}z\right|.\label{integ}
\end{equation}
记 $\lambda_i(z)=x_i dz/z$，便得到
\begin{equation}
\text{（$z$ 处的张力）}\ge|\lambda(z)|\qquad\text{其中}\ \lambda(z)=\lambda_1(z)-\lambda_2(z).
\end{equation}

% [Tutor Note: p.62 原文 "can be regarded as a string with variable string whose tension…" 中 "variable string" 应为 "variable tension"（上文即说 tension of the string）；译文按正意。]

M 理论中确实存在具有这些性质的高维对象。考虑形如
\begin{equation}
\mathbb{R}^{3,1}\times X\times\mathbb{R}^3.
\end{equation}
的十一维时空。M 理论具有称为 M5 膜的六维对象。将一个 M5 膜放在
\begin{equation}
\mathbb{R}^{3,1}\times\Sigma\times\{0\}
\end{equation}
上，其中 $\Sigma\subset X$ 是曲线，$0$ 是额外 $\bR^3$ 的原点。这给出一个四维理论。M 理论还具有称为 M2 膜的三维对象，它们的端点可以落在 M5 膜上。可以取一个 M2 膜位于
\begin{equation}
\mathbb{R}^{0,1}\times\text{圆盘}\times\{0\}
\end{equation}
其中 $\mathbb{R}^{0,1}\subset\mathbb{R}^{3,1}$ 是四维时空中粒子的世界线，而 $\text{圆盘}\subset X$ 的边界位于 $\Sigma$ 上，如图 \ref{fig:membrane-to-string} 所示。关于这一点的更多细节，读者可以从原始论文 \cite{Mikhailov:1997jv} 入手。

把系统视为定义在 $\mathbb{R}^{3,1}\times C$ 上的六维系统，考察这一中间情形也很有用。这个六维理论称为 $\SU(2)$ 型六维 $\cN{=}(2,0)$ 理论。

\subsubsection{例：纯 $\SU(2)$ 理论}
将这一高维思想具体应用到纯 $\SU(2)$ 理论的曲线 \eqref{curve_pure}。为便于查阅，重写曲线如下：
\begin{equation}
\Sigma:\qquad\Lambda^2z+\frac{\Lambda^2}z=x^2-u.
\end{equation}
将 $\Sigma$ 视为嵌入四维空间 $X$ 中。给定 $C$ 上的一点 $z$，求解上述二次方程便得到两个 $x$ 坐标，如图 \ref{fig:wboson_from_membrane} 左侧所示。将这两个解记为 $\pm x(z)$。随着点 $z$ 在 $C$ 上移动，它们形成曲线 $\Sigma$ 的两个叶片，见图 \ref{fig:wboson_from_membrane} 右侧。坐标 $x$ 总是用来描述 $C$ 上给出张力的一形式，因此始终将它乘以 $dz/z$ 会更方便，并可说两个叶片的坐标为 $\pm\lambda=x(z)dz/z$。从现在起采用这一约定。
\begin{figure}[h]
\[
\inc{wboson-as-membraneX}
\]
\caption{作为弦和悬置膜的 W 玻色子。\label{fig:wboson_from_membrane}}
\end{figure}

现在可以考虑在 $A$ 闭链上方、悬置于两个叶片之间的环状膜，见图 \ref{fig:wboson_from_membrane}。注意，作为 $C$ 上的一条弦，它的张力为 $2\lambda$，质量满足
\begin{equation}
M\ge|2\int_A\lambda|=|2a|.
\end{equation}
求解式 \eqref{curveBPSeq} 可以使张力最小，得到质量为
\begin{equation}
M=|2a|.
\end{equation}
的构型。注意，这恰好是 W 玻色子的正确质量；在我们的归一化下，它具有电荷 $n=2$，对应于 $\SU(2)$ 的三重态。还需注意，不可能得到质量为
\begin{equation}
M'=|a|,
\end{equation}
的膜，因为根本无法将膜悬置成只有一端位于 $A$ 闭链上方的情形。因此，相比于仅仅将曲线 $\Sigma$ 视为产生具有正确单值变换性质的全纯函数 $a(u)$、$a_D(u)$ 的辅助对象，这种高维推理具有更强的解释力。这一方法知道，这个系统中不存在电荷为 $n=1$ 的动力学粒子。
\begin{figure}[h]
\[
\inc{monopole-as-membraneX}
\]
\caption{作为弦和悬置膜的单极子。\label{fig:monopole_as_membrane}}
\end{figure}
接下来，可以考虑悬置于 $\Sigma$ 的叶片之间的圆盘状膜，使其端点位于 $C$ 的分支点 $z_+$、$z_-$ 上方，见图 \ref{fig:monopole_as_membrane}。通过与上面类似的推理，这个膜的质量为
\begin{equation}
M=2|\int_{z_-}^{z_+}\lambda|=|\int_B\lambda|=|a_D|.
\end{equation}
这是磁荷为 $m=1$ 的单极子的正确质量公式。从 $C$ 上张力可变的弦来看，需要注意，它对应于一条开弦，其端点位于张力 $2\lambda$ 变为零的点上。
\begin{figure}[h]
\[
\inc{Dyon-as-membrane}
\]
\caption{作为弦和悬置膜的双荷子。注意，它自动具有荷 $a_D+2a$，而\emph{不是} $a_D+a$。\label{fig:dyon_as_membrane}}
\end{figure}
还可以沿 $z$ 的相位方向绕行，将两个分支点 $z_\pm$ 连接起来，如图 \ref{fig:dyon_as_membrane} 所示。如图所示，这个膜在拓扑上是此前两个构型的和，其质量为
\begin{equation}
M=|2a+a_D|.
\end{equation}
这是电荷 $n$ 与磁荷 $m$ 满足 $(n,m)=(2,1)$ 的双荷子的正确质量公式。连接分支点时绕行 $n$ 次，就能得到质量为 $|2na+a_D|$ 的双荷子，其中 $n$ 为整数。同时也可看出，无论如何连接这些分支点，都无法得到质量为 $|(2n+1)a+a_D|$ 的双荷子，这与第 \ref{sec:monopole} 小节中的场论分析相容。

\subsection{六维理论的自对偶性}\index{SelfDuality@自对偶性|hyperpage}\index{SelfDualTwoForm@自对偶二形式|hyperpage}\index{DiracPairing@狄拉克配对|hyperpage}
现在我们发现，同一类对象，即 M 理论中的膜，或等价地六维 $\cN{=}(2,0)$ 理论中的弦，既能产生 W 玻色子这样的带电荷对象，也能产生单极子这样的带磁荷对象，见图 \ref{fig:wboson_from_membrane} 和图 \ref{fig:monopole_as_membrane}。为了把握这一性质，先回忆四维带荷粒子的基本特征，见图 \ref{fig:charged_objs}。
\begin{figure}[h]
\[
\inc{chargedthings}\qquad
\]
\caption{四维和六维中的带荷对象。\label{fig:charged_objs}}
\end{figure}
在一次量子化框架中，位于空间原点、沿时间方向 $t$ 延伸的电粒子，通过
\begin{equation}
\int_\text{世界线}A
\end{equation}
与电磁势耦合，从而产生非零电场 $F_{tr}\neq0$，其中
\begin{equation}
F=dA
\end{equation}
而 $r$ 为径向方向。世界线之外的运动方程为
\begin{equation}
dF=d{}\star F=0
\end{equation}
注意，在四维洛伦兹空间中，作用于二形式的 $\star^2=-1$。因此，不能施加条件 $\star F=F$。

考虑六维中由二形式 $B$ 描述的理论，弦通过下列项与之耦合：
\begin{equation}
\int_\text{世界面}B.
\end{equation}
设弦沿空间方向 $s$ 和时间方向 $t$ 延伸。这一构型产生非零电场 $G_{tsr}$，其中 $r$ 仍为径向方向。世界面之外的运动方程为
\begin{equation}
dG=d{}\star G=0
\end{equation}
这里 $\star$ 为六维 Hodge 星运算，定义为
\begin{equation}
(\star G)_{\mu\nu\rho}=\epsilon_{\mu\nu\rho\alpha\beta\gamma}G^{\alpha\beta\gamma}.
\end{equation}

在洛伦兹号差的六维空间中，作用于三形式的 $\star^2=1$，因此可以要求运动方程具有形式
\begin{equation}
dG=0,\quad G={}\star G.
\end{equation}
于是，沿 $t$ 和 $s$ 方向延伸的世界面，会同时具有非零电场 $G_{tsr}$ 和非零磁场 $G_{\theta\phi\psi}$。

% [Tutor Note: p.65 原文 "In six dimensions with Lorentizan signature" 中 "Lorentizan" 为 "Lorentzian" 之误；译文按正意。]

\begin{figure}[h]
\[
\inc{em-from-6d}
\]
\caption{从六维中的同一类对象产生电粒子与磁粒子。\label{fig:em-from-6d}}
\end{figure}
现在将理论置于坐标为 $x_{5,6}$ 的二维环面上，并考虑分别沿各个方向缠绕的弦，如图 \ref{fig:em-from-6d} 所示。将六维三形式场强记为 $G_{ABC}$，其中 $A,\ldots$ 是六维时空指标。考虑
\begin{equation}
F_{\mu\nu}:=G_{6\mu\nu},\qquad
F_{D\ \mu\nu}:=G_{5\mu\nu}.
\end{equation}
便可以提取出四维二形式。六维自对偶关系 $G={}\star_6G$ 转化为等式
\begin{equation}
F_D={}\star_4F.
\end{equation}
因此，六维中的单个自对偶二形式场产生一个 $\U(1)$ 场强。

缠绕 $x_6$ 的弦具有非零的 $G_{6tr}$ 和 $G_{5\theta\phi}$，因此具有非零的 $F_{tr}$。所以它在四维中成为一个电粒子。类似地，缠绕 $x_5$ 的弦具有非零的 $G_{5tr}$ 和 $G_{6\theta\phi}$。因而它具有非零的 $F_{\theta\phi}$，意味着它在四维中是一个磁粒子。
\begin{figure}[h]
\[
\inc{intersection}
\]
\caption{W 玻色子的膜与单极子的膜，其边界在两个点相交。\label{fig:intersection}}
\end{figure}
在纯 $\SU(2)$ 理论的具体情形中，W 玻色子和单极子由图 \ref{fig:intersection} 所示的膜产生。可以看到，W 玻色子的膜的边界与单极子的膜的边界在两个点相交。一般而言，只要计入取向所给出的符号，就可以通过这种交点计数得到它们作为四维时空中粒子的狄拉克配对。

\subsection{中间的五维杨--米尔斯理论及其边界条件}
\subsubsection{五维极大超对称杨--米尔斯理论}\index{FiveDimensionalYangMills@五维杨--米尔斯理论|hyperpage}\index{MonopoleString@单极子弦|hyperpage}
\begin{figure}[h]
\[
\inc{5dmsym-from-6dX}
\]
\caption{由六维得到的五维极大超对称杨--米尔斯理论。图中画出了一个 W 玻色子和一条单极子弦。\label{fig:5d-from-6d}}
\end{figure}
迄今为止，我们考虑的都是将六维理论置于某个二维空间上的情形，设其坐标为 $x_5$ 和 $x_6$。可以取一个极限，使 $x_5$ 方向远大于 $x_6$ 方向。此时，可以先沿 $x_6$ 方向紧致化，考虑一个中间的五维理论，见图 \ref{fig:5d-from-6d}。人们认为，这将给出规范群为 $\SU(2)$ 的五维极大超对称杨--米尔斯理论。

沿 $x_6$ 方向缠绕的弦产生一个有质量电粒子，而不沿 $x_6$ 方向缠绕的弦则成为一条有质量磁弦。这与 $\SU(2)$ 破缺为 $\U(1)$ 的五维杨--米尔斯理论的基本特征一致：首先，存在带电的有质量 W 玻色子；其次，只需规定场不依赖额外的第五个方向，就可以将四维规范理论的标准单极子解视为五维规范理论的解。此时，这些解应当解释为磁单极子弦。

\subsubsection{$\cN{=}4$ 超杨--米尔斯理论}\index{NFourTheory@$\cN{=}4$ 理论|hyperpage}\index{SenConjecture@Sen 猜想|hyperpage}
在 $x^5$ 方向施加周期性边界条件，就得到图 \ref{fig:n4sym-from-6d} 中的情形。我们正在将五维极大超对称杨--米尔斯理论紧致化到 $S^1$ 上，因此应当得到四维 $\cN{=}4$ 超杨--米尔斯理论。
\begin{figure}[h]
\[
\inc{n4sym-from-6dX}
\]
\caption{由六维得到的 $\cN{=}4$ 超杨--米尔斯理论。图中画出了一个 W 玻色子和一个单极子。\label{fig:n4sym-from-6d}}
\end{figure}
现在，\Gaiotto 曲线 $C$ 本身就是一个环面 $T^2$。令此 $T^2$ 的复结构为 $\tau$。\SeibergWitten\ 曲线 $\Sigma$ 由嵌入 $X$ 中的这个环面的两个平行副本组成，它们沿 $x$ 方向相隔 $2\lambda$，其中 $\lambda$ 现在为常数。

可以考虑 $C$ 中的闭链 $L_{n,m}$，它沿 $x_6$ 方向缠绕 $n$ 次，沿 $x_5$ 方向缠绕 $m$ 次。随后，考虑位于这一闭链上方的环状膜，就产生质量为
\begin{equation}
M_{n,m}=2|na+ma_D|
\end{equation}
的粒子，其中
\begin{equation}
a=\int_A\lambda,\quad
a_D=\int_B\lambda=\tau a.
\end{equation}
$(n,m)=(1,0)$ 的粒子是 W 玻色子，$(n,m)=(0,1)$ 的粒子是单极子。这一理论的特殊之处在于，单极子与 W 玻色子都来自环状膜。实际上，从六维观点看，区分环面的两个方向完全是任意的。因此，在交换 W 玻色子与单极子之后，耦合取某个值 $\tau=\tau_0$ 的理论，与取另一个值 $\tau=-1/\tau_0$ 的理论是相同的。

这些确实符合 $\cN{=}4$ 超对称 $\SU(2)$ 杨--米尔斯理论的性质。这一理论是共形的，并具有一个恰好边缘的耦合 $\tau$。在半经典区域中，单极子与 W 玻色子的质量比为 $|\tau|$。$\cN{=}4$ 超对称 $\SU(2)$ 杨--米尔斯理论具有四个伴随表示中的 Weyl 费米子。正如第 \ref{sec:monopole} 小节简要回顾过的，在这一情形下，对单极子解进行半经典量子化，会使单极子态组成一个有质量的 $\cN{=}4$ 矢量多重态。因而可以将其与同样处于有质量 $\cN{=}4$ 矢量多重态中的 W 玻色子交换。一般而言，我们预期对于任意互素整数对 $(n,m)$，都存在一个质量为 $|na+ma_D|$ 的有质量 $\cN{=}4$ 矢量多重态。它应当来自对荷为 $m$ 的单极子背景进行半经典量子化。这就是著名的 Sen 猜想 \cite{Sen:1994yi}。

\subsubsection{$\cN{=}2$ 纯 $\SU(2)$ 理论与 $N_f=1$ 理论}\index{HalfBPSBoundaryCondition@半 BPS 边界条件|hyperpage}
\begin{figure}[h]
\[
\inc{pure-and-nf1-via-5d}
\]
\caption{通过五维构造得到的纯 $\SU(2)$ 理论与 $N_f=1$ $\SU(2)$ 理论。\label{fig:pure_nf_1_from_5d}}
\end{figure}
$\cN{=}2$ 纯 $\SU(2)$ 理论的曲线
\begin{equation}
\frac{\Lambda^2}z+\Lambda^2z=x^2-u
\end{equation}
以及含有一个味的 $\cN{=}2$ $\SU(2)$ 理论的曲线
\begin{equation}
\frac{2\Lambda(x-\mu)}z+\Lambda^2z=x^2-u
\end{equation}
都可以赋予类似的解释。要点是取 $x_5=\log|z|$、$x_6=\Arg z$，并先沿 $x_6$ 方向紧致化，见图 \ref{fig:pure_nf_1_from_5d}。

先考虑纯理论。左端的 $\Lambda^2/z$ 项应当视为使第五个方向 $x_5$ “终止”的边界条件，尽管 $x_5=\log|z|$ 在形式上延伸至 $-\infty$。五维理论的内部（bulk）具有极大超对称性。得到的四维理论具有 $\cN{=}2$ 超对称性，因此边界破缺了一半超对称性，除此之外并未产生太多其他影响。保持原来一半超对称性的边界条件称为半 BPS 边界条件。由此可见，$\Lambda^2/z$ 项代表五维理论的一个半 BPS 边界条件。

$\Lambda^2z$ 项由翻转 $x_5\leftrightarrow-x_5$ 得到，因此应当代表同样的边界条件。最终，我们看到，这一系统是将极大超对称 $\SU(2)$ 杨--米尔斯理论紧致化在一条线段上，并由两个破缺一半超对称性的边界条件使其终止，从而实现四维纯 $\SU(2)$ 杨--米尔斯理论。

接下来考虑单味理论。$\Lambda^2z$ 项与纯理论相同，因此应给出同样的半 BPS 边界条件。$z\sim0$ 处的边界条件则不同：现在有一个形如 $\Lambda(x-\mu)/z$ 的项。这应当意味着，边界上存在一个质量为 $\mu$、处于 $\SU(2)$ 二重态中的超多重态，并与体内的五维规范多重态耦合。

\subsubsection{含有 $N_f=2,3,4$ 个味的 $\SU(2)$ 理论}\label{sec:su2nfpre}
由这一解释，很容易得到含有 $N_f=2,3,4$ 个味的 $\SU(2)$ 理论的六维实现，也就是含有 $N_f=2,3,4$ 个二重态表示超多重态的理论。用 $\cN{=}1$ 手征多重态表示时，场为 $(Q_i^a,\tilde Q^i_a)$，其中 $a=1,2$、$i=1,\ldots,N_f$，超势为
\begin{equation}
\sum_i\left(Q_i\Phi\tilde Q^i+\mu_iQ_i\tilde Q^i\right)
\end{equation}
其中 $\mu_i$ 为质量项。

从 $N_f=2$ 理论开始。我们已经知道，如何在 $z=0$ 一侧的边界上引入一个二重态超多重态。要在 $z=\infty$ 一侧做同样的事情，只需作变量替换 $z\leftrightarrow1/z$。最终得到图 \ref{fig:nf_2_from_6d} 左侧所示的设置，其曲线为
\begin{equation}
\frac{2\Lambda(x-\mu_1)}z+2\Lambda(x-\mu_2)z=x^2-u\label{nf_2_curve_pre}
\end{equation}
并配以 Seiberg-Witten 微分 $\lambda=xdz/z$。

% [Tutor Note: p.70 原文 "we just a change of variables $z\leftrightarrow 1/z$" 脱动词（应为 "we just make a change of variables"）；译文按正意。]
\begin{figure}[h]
\[
\inc{su2nf2-from-6dXX}
\]
\caption{$N_f=2$ 理论。\label{fig:nf_2_from_6d}}
\end{figure}

引入另一个变量 $z'=(x-\mu_2)z/(2\Lambda)$，同一条曲线可以重写为
\begin{equation}
\frac{(x-\mu_1)(x-\mu_2)}{z'}+4\Lambda^2z'=x^2-u.\label{nf_2_curve_pre2}
\end{equation}
但现在可以取 $x_5'=\log|z'|$、$x_6'=\Arg z'$，先约化为由 $z'$ 参数化的 $C'$ 上的理论，再约化为由 $|z'|$ 参数化的线段上的五维理论。在这种解释中，$z'=\infty$ 一侧的边界条件与纯 $\SU(2)$ 情形相同。因此，$z'=0$ 一侧由 ${(x-\mu_1)(x-\mu_2)}/z$ 项给出的边界条件，应当是使边界上存在两个 $\SU(2)$ 二重态超多重态的半 BPS 条件。

只有给出描述可变张力的一形式 $\lambda$，系统的描述才是完整的。在式 \eqref{nf_2_curve_pre} 中，它为 $2\pi i\lambda=xdz/z$；在式 \eqref{nf_2_curve_pre2} 中，它为 $2\pi i\lambda'=xdz'/z'$。二者都由积分 $dx\wedge d\log z=dx\wedge d\log z'$ 得到，见式 \eqref{integ}。不过，这两个微分形式并不完全相等：
\begin{equation}
\lambda'-\lambda=\frac{1}{2\pi i}xd\log\frac{z'}z=\frac{1}{2\pi i}xd\log(x-\mu_2).\label{diff}
\end{equation}
其差与 $u$ 无关，在 $x=\mu_2$ 处具有非零留数 $\mu_2$。这意味着，给定 \SeibergWitten\ 曲线 $\Sigma$ 上的一个闭链 $L$，有
\begin{equation}
\oint_L\lambda'-\oint_L\lambda=k\mu_2
\end{equation}
其中 $k$ 为整数。回忆一下，BPS 质量公式由展开
\begin{equation}
\oint_L\lambda=na+ma_D+f_1\mu_1+f_2\mu_2
\end{equation}
决定，其中 $f_{1,2}$ 是味荷，见式 \eqref{Zlin}。因此，在两个 Seiberg-Witten 微分 $\lambda$ 和 $\lambda'$ 之间作选择，会影响味荷 $f_2$ 与闭链 $L$ 之间的对应关系，但几乎不影响其他内容。一般而言，如果对 Seiberg-Witten 微分所作的改变，是加上一个与 $u$ 无关、且留数为超多重态质量的整数线性组合的微分形式，那么这种改变是安全的。后面我们还会反复遇到这种情形。

% [Tutor Note: p.71 原文 "its non-zero residue is at μ2 at x=μ2" 表述含混（应读作：唯一的非零留数位于 $x=\mu_2$ 处、其值为 $\mu_2$）；译文按此理解。]
\begin{figure}[h]
\[
\inc{su2nf3nf4-from-6dX}
\]
\caption{$N_f=3$ 理论与 $N_f=4$ 理论。\label{fig:nf_3_4_from_6d}}
\end{figure}
既然已经得到了代表两个二重态超多重态的边界条件，就很容易猜出 $N_f=3$ 和 $N_f=4$ 理论的曲线。只需将已经找到的各种边界条件组合起来，如图 \ref{fig:nf_3_4_from_6d} 所示。对 $N_f=3$ 理论，得到
\begin{equation}
\frac{(x-\mu_1)(x-\mu_2)}z+2\Lambda(x-\mu_3)z=x^2-u,\label{nf_3_curve_pre}
\end{equation}
对 $N_f=4$ 理论，得到
\begin{equation}
f\cdot\frac{(x-\mu_1)(x-\mu_2)}z+f'\cdot(x-\mu_3)(x-\mu_4)z=x^2-u\label{nf_4_curve_pre}
\end{equation}
其中引入了复数 $f$ 和 $f'$。其中一个可以通过重新缩放 $z$ 消去。

接下来的任务，是检验通过六维构造得到的这些曲线，是否具有描述相应四维理论所需的正确性质。在继续之前，我们还需要进一步了解一般 $\cN{=}2$ 理论的希格斯支。
